% Redrawn 2026-10-15; main labels 8 pt, notes 7.5 pt; actual fonts embedded in PDF.
% export-reference-width-mm: 112
% The return edge is available only under an online data-collection protocol.
\begin{tikzpicture}[
  x=1mm,y=1mm,font=\figurefont\fontsize{8}{10}\selectfont,
  box/.style={rounded corners=1.8pt,draw=FigureStroke,line width=.6pt,fill=FigureBlueFill,
    minimum width=36mm,minimum height=13mm,align=center,inner sep=1.5mm},
  flow/.style={-{Latex[length=2mm]},draw=FigureStroke,line width=.6pt,
    rounded corners=1.8pt},
  online/.style={flow,dashed,draw=BookAmber},
  tag/.style={fill=white,inner sep=.7mm}
]
  \node[box] (data) at (23,49) {真实交互训练记录\\观测、动作、奖励与终止};
  \node[box,fill=FigureGreenFill] (model) at (82,49) {世界模型\\动作条件下预测后果};
  \node[box,fill=FigureYellowFill] (solve) at (82,13) {模型内求解\\训练策略或局部规划};
  \node[box] (real) at (23,13) {真实执行与评价\\回报、误差与失败};
  \draw[flow] (data) -- node[tag,above] {拟合} (model);
  \draw[flow] (model) -- node[tag,right] {滚动预测} (solve);
  \draw[flow] (solve) -- node[tag,above] {策略或动作} (real);
  \draw[online] (real) -- node[tag,left,align=right] {新交互记录\\校准模型} (data);
  \node[anchor=north,align=center,text=BookMuted,font=\FigureNoteFont\fontsize{7.5}{9.5}\selectfont] at (52,-1)
    {虚线仅在允许追加交互训练记录时存在；\\纯离线方法不能把测试记录回流训练。};
\end{tikzpicture}
